% =========================================================
% Tactics Catalog
% =========================================================
\paragraph{Catalog of Algebraic Tactics.}

The following catalog lists the specific algebraic moves available once
a proof is in progress and the architecture has been chosen. Each entry
states the tactic, its precondition (what must already be established
or assumed), and its effect (what it achieves in the proof).

\begin{center}
\renewcommand{\arraystretch}{1.5}
\begin{tabular}{p{0.20\textwidth} p{0.35\textwidth} p{0.35\textwidth}}
\toprule
\textbf{Tactic} & \textbf{Precondition} & \textbf{Effect} \\
\midrule

\textbf{Multiply by inverse} &
  $a \neq 0$ in a field; or $a \in G$ in a group &
  Multiply both sides by $a^{-1}$; use associativity to collapse
  $a^{-1}a = e$; isolates the target variable. \\

\textbf{Cancellation} &
  $ac = bc$ (or $ca = cb$) in a group or integral domain &
  Conclude $a = b$. In a group: multiply by $c^{-1}$.
  In an integral domain: cite no-zero-divisors. \\

\textbf{Distributivity bridge} &
  Multiplication by $0$, $-1$, or a negative element &
  Rewrite the product as a sum using distributivity, then use
  additive group structure to collapse. Crosses from
  $\times$ to $+$. \\

\textbf{Add zero} &
  Need to introduce a term without changing the expression &
  Write $a = a + 0 = a + (b + (-b))$; regroup to expose
  the needed structure. \\

\textbf{Satisfy-and-cite} &
  A uniqueness theorem for some named object $b$ is in scope &
  Show $a$ satisfies $b$'s defining property; cite uniqueness;
  conclude $a = b$. Avoids assume-two-and-compare entirely. \\

\textbf{Double negation} &
  $-(-a)$ appears or is the goal &
  $-a$ is the unique element $x$ with $a + x = 0$;
  $-(-a)$ satisfies this; by uniqueness, $-(-a) = a$. \\

\textbf{Absorb into axiom} &
  A subexpression matches the LHS of an axiom &
  Rewrite using the axiom. Tag as \tagTA. \\

\bottomrule
\end{tabular}
\end{center}

\begin{remark}[The distributivity bridge in detail]
The distributivity bridge is the central tactic for ring and field
proofs involving zero and negatives. Every instance of ``$a \cdot 0 = 0$''
and ``$a(-b) = -(ab)$'' uses it. Once you have seen the pattern, it is
immediately recognizable.

The pattern always has four steps:
\begin{enumerate}[itemsep=2pt]
  \item Introduce $a \cdot 0$ or $a \cdot (-b)$ into the argument.
  \item Expand: write $0 = x + (-x)$ or $-b = b + (-b) + (-b)$ using
    the additive inverse definition.
  \item Apply distributivity: $a(x + y) = ax + ay$.
  \item Use uniqueness of the additive inverse to identify the result
    with $-(ab)$ or $0$.
\end{enumerate}

The reason this requires distributivity as a bridge is that the
additive structure (with its identity $0$ and its inverses $-a$) and
the multiplicative structure (with its products $ab$) are separate. The
only connection between them is distributivity. To get from the multiplicative
side ($a \cdot 0$) to the additive side ($0 = 0 + 0$), you must cross
the bridge.
\end{remark}

\begin{remark}[Cancellation: two theorems that must not be conflated]
``Cancellation'' names two distinct results with different preconditions.

\medskip
\noindent\textbf{Group cancellation.} In any group, $ac = bc \Rightarrow a = b$.
Proved by multiplying both sides on the right by $c^{-1}$, which always
exists in a group. The hypothesis $c \neq 0$ is not needed because
inverses always exist in a group.

\medskip
\noindent\textbf{Domain cancellation (or integral domain cancellation).}
In an integral domain, $ac = bc$ and $c \neq 0$ imply $a = b$.
Proved by rewriting as $(a - b)c = 0$ and invoking the no-zero-divisors
property. The hypothesis $c \neq 0$ is essential: without it, cancellation
can fail (e.g., in $\mathbb{Z}_6$, $2 \cdot 3 = 4 \cdot 3$ but $2 \neq 4$).

\medskip
\noindent
The error to avoid: using group cancellation in an integral domain
argument (it works but uses stronger machinery than needed and may not
be available), or trying to use domain cancellation without verifying
$c \neq 0$.
\end{remark}

\begin{remark}[Add zero: the simplest but most overlooked tactic]
The add-zero tactic sounds trivial but appears in almost every algebraic
proof where an intermediate term needs to be introduced without
changing the value of an expression. The template is:
\[
a = a + 0 = a + (b + (-b)) = (a + b) + (-b).
\]
This exposes $a + b$ as a subexpression of $a$, which can then be
worked with. It is the algebraic equivalent of adding and subtracting
the same term in a real analysis inequality argument.

When you need a term to appear in an expression and cannot see how to
introduce it, ask: can I write some zero as $x + (-x)$ and distribute?
\end{remark}
